\section*{Sets and Indices}

\begin{itemize}
    \item $F$: set of all foods, indexed by $f$.
    \item $V \subseteq F$: set of vegetables.
    \item $P \subseteq F$: set of proteins.
\end{itemize}

For this instance,
\[
V=\{\text{okra},\text{carrots},\text{celery},\text{cabbage}\},\qquad
P=\{\text{salmon},\text{beef},\text{pork}\},
\]
and
\[
F = V \cup P.
\]

\section*{Parameters}

\begin{itemize}
    \item $b$: meal budget in dollars (code: \texttt{budget}); here $b=15$.
    \item $T$: required total food intake in grams (code: \texttt{food\_intake}); here $T=600$.
    \item $\phi_f$: fiber content per 100g of food $f$ (code: \texttt{fiber[f]}). Missing values in the data are filled with $0$, so the protein items have $\phi_f=0$.
    \item $c_f$: price per 100g of food $f$ in dollars (code: \texttt{price[f]}).
    \item $M$: big-$M$ constant used in the linking constraints (code: \texttt{M}), defined as
    \[
    M = T/100.
    \]
\end{itemize}

Numerically for this instance,
\[
T/100 = 6,
\]
and
\[
\begin{aligned}
&\phi_{\text{okra}}=3.2,\ \phi_{\text{carrots}}=2.7,\ \phi_{\text{celery}}=1.6,\ \phi_{\text{cabbage}}=2.0,\\
&\phi_{\text{salmon}}=0,\ \phi_{\text{beef}}=0,\ \phi_{\text{pork}}=0,
\end{aligned}
\]
\[
\begin{aligned}
&c_{\text{okra}}=2.6,\ c_{\text{carrots}}=1.2,\ c_{\text{celery}}=1.6,\ c_{\text{cabbage}}=2.3,\\
&c_{\text{salmon}}=4.0,\ c_{\text{beef}}=3.6,\ c_{\text{pork}}=1.8.
\end{aligned}
\]

\section*{Decision Variables}

\begin{itemize}
    \item $x_f \ge 0$ (code: \texttt{x[f]}): amount of food $f$ selected, measured in units of 100g.
    \item $y_f \in \{0,1\}$ (code: \texttt{y[f]}): selection indicator for food $f$.
\end{itemize}

\section*{Objective}

Maximize total fiber intake:
\[
\max \sum_{f \in F} \phi_f x_f.
\]

\section*{Constraints}

Total intake must equal the required food amount:
\[
\sum_{f \in F} x_f = T/100.
\]

Budget limit:
\[
\sum_{f \in F} c_f x_f \le b.
\]

Exactly one protein source must be selected:
\[
\sum_{f \in P} y_f = 1.
\]

At least two vegetable types must be selected:
\[
\sum_{f \in V} y_f \ge 2.
\]

Linking constraints with minimum inclusion amount:
\[
x_f \ge 0.1\, y_f \qquad \forall f \in F,
\]
\[
x_f \le M\, y_f \qquad \forall f \in F.
\]

\section*{Notes on Implementation}

\begin{itemize}
    \item The model is implemented as a mixed-integer linear program and solved with Gurobi through PuLP's compatible interface (\texttt{pulp.GUROBI\_CMD}).
    \item The quantity variables are measured in 100g units, so the 600g intake requirement becomes $6$ in the model.
    \item The lower linking bound $x_f \ge 0.1\,y_f$ enforces that any selected food must appear in an amount of at least $0.1$ units, i.e., 10g.
    \item The upper linking bound uses $M=T/100=6$, which is exactly the value used in the code.
    \item Fiber values missing in the CSV are replaced by zero before optimization; thus all protein items contribute zero fiber in the objective.
    \item The binary variables are defined for all foods, not only for proteins and vegetables, and the cardinality constraints are applied separately to the subsets $P$ and $V$.
\end{itemize}
